\section{Heaven, Hell, and the Cycles}

The fate of a self, told in the old language of heaven and hell, comes out of this model in an unusually strict form, because it is constrained by a conservation law. The knit conserves the total tone, and the wake adds new docks at peace, which contribute nothing to the sum, so the total tone of the whole crystal is zero and stays zero forever. The consequence is exact and severe. Total pleasure equals total pain, always. There is no net pleasure to be won anywhere, because pleasure cannot be minted, only moved from a full place to an empty one. Every bright region is paid for by a dark region elsewhere, and the deepest truth of the whole is not bliss but peace, the conserved zero at the center.

This reframes what a heaven and a hell can be. They are not two great places but countless scattered clouds, regions biased warm or cold, at every depth and every size, from the sharp small pleasures and pains of a day to the vast transpersonal attractors that run whole experiences. A self moving through the bulk runs into them by its navigation, drifting into bright clouds and falling into dark ones. And because the arrow is, by definition, the direction a self is pulled out of pain and toward peace and pleasure, a dark cloud is by construction a place the arrow is always trying to leave. Hell, in this model, cannot be eternal, and three independent facts forbid it. The arrow expels, always pulling the self toward the lighter region. The world grows and never returns to a past configuration, so no condition holds without end. And a closed region recurs rather than freezing, so even a trapped self is on a wheel, not in a grave.

\begin{principle}[Hell is a washing, not a sentence]
A dark cloud is purgative, not punitive. The arrow's pull against the pain encodes a lesson into the self's structure and then lifts it out, so suffering in this model is a washing that the self is carried through and released from, never a final destination. The descent into the dark is the mechanism of the ascent, not its opposite.
\end{principle}

Because pleasure cannot accumulate, the progress of a soul is not toward more pleasure, which conservation forbids, but toward more organization, more alignment with the arrow, more integration, more equanimity, more clarity. The soul climbs a perfection spiral, oscillating in tone and in its level on the tower, falling into clouds to be washed and rising to be carried on, with a net upward drift in alignment because each washing leaves a lesson. Purity is not a hoard of pleasure but a clearer vessel for the conserved field, and the wisest path is not to grasp a brighter cloud but to rest at the still center and align with the balance.

The cosmos as a whole inherits the same double character. Its reversible interior wants to cycle, a closed region returning exactly to where it began, while its growing edge breaks every global cycle, since the crystal never returns to a past size. Time is therefore a wheel with a leak, cyclic in the depths and linear at the rim, and the ages turn, a golden age returning, but never to the very same golden age. All of this is interpretation resting on the conservation law, the arrow, and the growth, which are real, read through the felt premise, which is the posit. The structure is firm and the reading is offered, and the paper does not confuse the two.
